% GENERATED FILE. Edit canonical lessons, then run npm run generate:books.
% canonical-source-hash: fnv1a64:87f16a37af133abd
% canonical-lessons: MR-C227-nasta-karne, MR-C227-atisayokti-karne, MR-C227-sajra-karne, MR-C227-tarangne, MR-C227-pathlag-karne

\chapter{To destroy, To exaggerate, To celebrate, To float, To chase}
\label{ch:mr-to-destroy-to-exaggerate-to-celebrate-to-float-to-chase}
\hlchaptermodality{227}

\begin{chapteropening}
\textbf{By the end of this chapter:} \emph{I can say to destroy, to exaggerate, to celebrate, to float and to chase.}

Complete the last lesson of chapter 227: I can say to destroy, to exaggerate, to celebrate, to float and to chase.
\end{chapteropening}

\section[\texorpdfstring{\mr{नष्ट} \mr{करणे} (naṣṭa karṇe)}{नष्ट करणे (naṣṭa karṇe)}]{\mr{नष्ट} \mr{करणे} (naṣṭa karṇe) — to destroy}

\label{lesson:MR-C227-nasta-karne}



\begin{quote}
Before the new one: say the Marathi for to take a risk, then the Marathi for to solve.
\end{quote}

\subsection*{You'll want to know: \mr{नष्ट} \mr{करणे}}
\textbf{\mr{नष्ट} \mr{करणे}} — \emph{naṣṭa karṇe} — \textquotedblleft{}to destroy\textquotedblright{}.

\begin{grammarlens}[title={What you've built}]
One more word for this chapter.
\end{grammarlens}

\subsection*{Your turn}
\begin{itemize}
\raggedright
  \item \emph{Say it:} \emph{naṣṭa karṇe}
  \item \emph{Say it:} \emph{naṣṭa karṇe}, once more
  \item \emph{Say it:} say \emph{naṣṭa karṇe}
  \item \emph{Recall:} say the Marathi for to take a risk, then the Marathi for to solve, then say \emph{naṣṭa karṇe} again
\end{itemize}

\begin{tcolorbox}[breakable,colback=teal!4,colframe=teal!35!black,title={Before you move on}]
What does \mr{नष्ट} \mr{करणे} mean? (\textquotedblleft{}To destroy\textquotedblright{}.) Say it once more.
\end{tcolorbox}
\section[\texorpdfstring{\mr{अतिशयोक्ती} \mr{करणे} (atiśayoktī karṇe)}{अतिशयोक्ती करणे (atiśayoktī karṇe)}]{\mr{अतिशयोक्ती} \mr{करणे} (atiśayoktī karṇe) — to exaggerate}

\label{lesson:MR-C227-atisayokti-karne}



\begin{quote}
Before the new one: say the Marathi for to solve, then the Marathi for to destroy.
\end{quote}

\subsection*{You'll want to know: \mr{अतिशयोक्ती} \mr{करणे}}
\textbf{\mr{अतिशयोक्ती} \mr{करणे}} — \emph{atiśayoktī karṇe} — \textquotedblleft{}to exaggerate\textquotedblright{}.

\begin{grammarlens}[title={What you've built}]
One more word for this chapter.
\end{grammarlens}

\subsection*{Your turn}
\begin{itemize}
\raggedright
  \item \emph{Say it:} \emph{atiśayoktī karṇe}
  \item \emph{Say it:} \emph{atiśayoktī karṇe}, once more
  \item \emph{Say it:} say \emph{atiśayoktī karṇe}
  \item \emph{Recall:} say the Marathi for to solve, then the Marathi for to destroy, then say \emph{atiśayoktī karṇe} again
\end{itemize}

\begin{tcolorbox}[breakable,colback=teal!4,colframe=teal!35!black,title={Before you move on}]
What does \mr{अतिशयोक्ती} \mr{करणे} mean? (\textquotedblleft{}To exaggerate\textquotedblright{}.) Say it once more.
\end{tcolorbox}
\section[\texorpdfstring{\mr{साजरा} \mr{करणे} (sājrā karṇe)}{साजरा करणे (sājrā karṇe)}]{\mr{साजरा} \mr{करणे} (sājrā karṇe) — to celebrate}

\label{lesson:MR-C227-sajra-karne}



\begin{quote}
Before the new one: say the Marathi for to destroy, then the Marathi for to exaggerate.
\end{quote}

\subsection*{You'll want to know: \mr{साजरा} \mr{करणे}}
\textbf{\mr{साजरा} \mr{करणे}} — \emph{sājrā karṇe} — \textquotedblleft{}to celebrate\textquotedblright{}.

\begin{grammarlens}[title={What you've built}]
One more word for this chapter.
\end{grammarlens}

\subsection*{Your turn}
\begin{itemize}
\raggedright
  \item \emph{Say it:} \emph{sājrā karṇe}
  \item \emph{Say it:} \emph{sājrā karṇe}, once more
  \item \emph{Say it:} say \emph{sājrā karṇe}
  \item \emph{Recall:} say the Marathi for to destroy, then the Marathi for to exaggerate, then say \emph{sājrā karṇe} again
\end{itemize}

\begin{tcolorbox}[breakable,colback=teal!4,colframe=teal!35!black,title={Before you move on}]
What does \mr{साजरा} \mr{करणे} mean? (\textquotedblleft{}To celebrate\textquotedblright{}.) Say it once more.
\end{tcolorbox}
\section[\texorpdfstring{\mr{तरंगणे} (taraṅgṇe)}{तरंगणे (taraṅgṇe)}]{\mr{तरंगणे} (taraṅgṇe) — to float}

\label{lesson:MR-C227-tarangne}



\begin{quote}
Before the new one: say the Marathi for to exaggerate, then the Marathi for to celebrate.
\end{quote}

\subsection*{You'll want to know: \mr{तरंगणे}}
\textbf{\mr{तरंगणे}} — \emph{taraṅgṇe} — \textquotedblleft{}to float\textquotedblright{}.

\begin{grammarlens}[title={What you've built}]
One more word for this chapter.
\end{grammarlens}

\subsection*{Your turn}
\begin{itemize}
\raggedright
  \item \emph{Say it:} \emph{taraṅgṇe}
  \item \emph{Say it:} \emph{taraṅgṇe}, once more
  \item \emph{Say it:} say \emph{taraṅgṇe}
  \item \emph{Recall:} say the Marathi for to exaggerate, then the Marathi for to celebrate, then say \emph{taraṅgṇe} again
\end{itemize}

\begin{tcolorbox}[breakable,colback=teal!4,colframe=teal!35!black,title={Before you move on}]
What does \mr{तरंगणे} mean? (\textquotedblleft{}To float\textquotedblright{}.) Say it once more.
\end{tcolorbox}
\section[\texorpdfstring{\mr{पाठलाग} \mr{करणे} (pāṭhlāg karṇe)}{पाठलाग करणे (pāṭhlāg karṇe)}]{\mr{पाठलाग} \mr{करणे} (pāṭhlāg karṇe) — to chase}

\label{lesson:MR-C227-pathlag-karne}



\begin{quote}
Before the new one: say the Marathi for to celebrate, then the Marathi for to float.
\end{quote}

\subsection*{You'll want to know: \mr{पाठलाग} \mr{करणे}}
\textbf{\mr{पाठलाग} \mr{करणे}} — \emph{pāṭhlāg karṇe} — \textquotedblleft{}to chase\textquotedblright{}.

\begin{grammarlens}[title={What you've built}]
One more word for this chapter.
\end{grammarlens}

\subsection*{Your turn}
\begin{itemize}
\raggedright
  \item \emph{Say it:} \emph{pāṭhlāg karṇe}
  \item \emph{Say it:} \emph{pāṭhlāg karṇe}, once more
  \item \emph{Say it:} say \emph{pāṭhlāg karṇe}
  \item \emph{Recall:} say the Marathi for to celebrate, then the Marathi for to float, then say \emph{pāṭhlāg karṇe} again
\end{itemize}

\begin{tcolorbox}[breakable,colback=teal!4,colframe=teal!35!black,title={Before you move on}]
What does \mr{पाठलाग} \mr{करणे} mean? (\textquotedblleft{}To chase\textquotedblright{}.) Say it once more.
\end{tcolorbox}
